% ── Záver ─────────────────────────────────────────────────────
\chapter*{Záver}
\addcontentsline{toc}{chapter}{Záver}

        Cieľom tejto práce bolo navrhnúť a~implementovať webový systém na správu
projektov s~integrovaným AI~chatom nad dokumentmi. Výsledná aplikácia ProMat napĺňa
všetky stanovené požiadavky a~predstavuje funkčný nástroj pre tímy, ktoré potrebujú
centralizovať riadenie projektov, úloh a~dokumentácie na jednom mieste.

        V~teoretickej časti sme analyzovali problematiku roztrieštenej práce s~nástrojmi,
porovnali existujúce riešenia (napríklad Trello) a~identifikovali ich kľúčové
obmedzenia, predovšetkým chýbajúcu inteligentnú prácu s~priloženými dokumentmi.
Na konci tejto časti bol popísaný technologický zásobník zvoleného riešenia.

        V~praktickej časti sme opísali architektonický návrh systému, 
autentifikáciu, správu projektov a~úloh vrátane Kanban tabuľe,
nahrávanie a~zdieľanie dokumentov, AI~chat nad ich obsahom a~informačný dashboard.

        Najvýznamnejším prínosom projektu je integrácia lokálneho LLM servera Ollama
s~modelom Gemma~4, vďaka ktorej môžu tímy využívať výhody umelej inteligencie pri
práci s~citlivými dokumentmi bez nutnosti odosielania dát na externé cloudové servery.
Toto riešenie je vhodné najmä pre organizácie s~prísnejšími bezpečnostnými alebo
legislatívnymi požiadavkami na ochranu dát.

        Predkladaný projekt preukázal, že prepojenie projektového manažmentu
s~inteligentnou prácou s~dokumentmi je technicky realizovateľné a~prakticky
prínosné riešenie, ktoré môže výrazne zvýšiť efektivitu tímov v~rôznych typoch
organizácií.

        Pri tvorbe aplikácie sa zároveň potvrdilo, že moderné open-source technológie
umožňujú vytvoriť robustné riešenie aj bez potreby komerčných platforiem. Kombinácia
frameworku Flask, relačnej databázy a~lokálneho LLM servera poskytla dostatočný výkon,
modularitu a~možnosť ďalšieho rozširovania podľa potrieb používateľov.

        Priestor pre budúci rozvoj vidíme najmä v~podpore integrácie s~externými systémami a~vo
vylepšení AI~chatu o~presnejšie citovanie zdrojových pasáží z~dokumentov.

        Na základe dosiahnutých výsledkov možno konštatovať, že stanovený cieľ práce bol
splnený v~plnom rozsahu. Aplikácia ProMat predstavuje prakticky použiteľný základ,
na ktorom je možné ďalej budovať komplexný interný nástroj pre projektové riadenie
a~efektívnu prácu s~dokumentmi podporenú umelou inteligenciou.